\begin{lemma}
\label{lemma-pseudo-coherent-descends-fppf}
Let $\{f_i : X_i \to X\}$ be an fppf covering of schemes. Let
$E \in D(\mathcal{O}_X)$. Let $m \in \mathbf{Z}$.
Then $E$ is $m$-pseudo-coherent if and only if each
$Lf_i^*E$ is $m$-pseudo-coherent.
\end{lemma}

\begin{proof}
Pullback always preserves $m$-pseudo-coherence, see
Cohomology, Lemma \ref{cohomology-lemma-pseudo-coherent-pullback}.
Conversely, assume that $Lf_i^*E$ is $m$-pseudo-coherent for all $i$.
Let $U \subset X$ be an affine open. It suffices to prove that
$E|_U$ is $m$-pseudo-coherent. Since $\{f_i : X_i \to X\}$ is an
fppf covering, we can find finitely many affine open $V_j \subset X_{a(j)}$
such that $f_{a(j)}(V_j) \subset U$ and $U = \bigcup f_{a(j)}(V_j)$.
Set $V = \coprod V_i$.
Thus we may replace $X$ by $U$ and $\{f_i : X_i \to X\}$ by
$\{V \to U\}$ and assume that $X$ is affine and our covering
is given by a single surjective flat morphism $\{f : Y \to X\}$
of finite presentation.

\medskip\noindent
Since $f$ is flat the derived functor $Lf^*$ is just given by $f^*$ and $f^*$
is exact. Hence $H^i(Lf^*E) = f^*H^i(E)$. Since $Lf^*E$ is $m$-pseudo-coherent,
we see that $Lf^*E \in D^-(\mathcal{O}_Y)$. Since $f$ is surjective and flat,
we see that $E \in D^-(\mathcal{O}_X)$. Let $i \in \mathbf{Z}$ be the largest
integer such that $H^i(E)$ is nonzero. If $i < m$, then we are done. Otherwise,
$f^*H^i(E)$ is a finite type $\mathcal{O}_Y$-module by
Cohomology, Lemma \ref{cohomology-lemma-finite-cohomology}.
Then by Descent, Lemma \ref{descent-lemma-finite-type-descends-fppf}
the $\mathcal{O}_X$-module $H^i(E)$ is of finite type.
Thus, after replacing $X$ by the members of a finite affine open covering,
we may assume there exists a map
$$
\alpha : \mathcal{O}_X^{\oplus n}[-i] \longrightarrow E
$$
such that $H^i(\alpha)$ is a surjection. Let $C$ be the cone of $\alpha$
in $D(\mathcal{O}_X)$. Pulling back to $Y$ and using
Cohomology, Lemma \ref{cohomology-lemma-cone-pseudo-coherent}
we find that $Lf^*C$ is $m$-pseudo-coherent. Moreover $H^j(C) = 0$
for $j \geq i$. Thus by induction on $i$ we see that $C$ is
$m$-pseudo-coherent. Using
Cohomology, Lemma \ref{cohomology-lemma-cone-pseudo-coherent}
again we conclude.
\end{proof}